% Part: sets-functions-relations
% Chapter: arithmetization
% Section: From N to Z

\documentclass[../../../include/open-logic-section]{subfiles}


\begin{document}

\olfileid{sfr}{arith}{int}
\olsection{$\Nat$ توں $\Int$ ول}

ایتھے دو بنیادی گلّاں نیں:
	\begin{enumerate}
		\item ہر صحیح عدد نوں $n - m$ دی صورت وچ لکھیا جا سکدا اے، جتھے $n, m \in\Nat$۔
		\item اظہار $n - m$ وچ کوڈ کیتی معلومات نوں ترتیب وار جوڑے $\tuple{n, m}$ نال وی اوہو جیہے طور تے کوڈ کیتا جا سکدا اے۔
	\end{enumerate}
اسی پہلے ای جان دے آں کہ قدرتی عدداں دیاں ترتیب وار جوڑیاں
$\Nat^2$ دے !!{element}s نیں۔ تے اسی فرض کر رہے آں کہ ساڈے نوں
$\Nat$ دی سمجھ اے۔ سو ایہناں دوہاں گلّاں دے آسرے اک سادی تجویز ایہ
اے: \emph{صحیح عدداں نوں قدرتی عدداں دیاں ترتیب وار جوڑیاں سمجھ لئیے}۔

اصل وچ ایہ تجویز حدوں ودھ سادی اے۔ ظاہر اے کہ اسی ایہ چاہنے آں کہ
$0- 2 = 4 - 6$ ہووے۔ پر صاف اے کہ $\tuple{0, 2 }\neq \tuple{4, 6}$۔
سو اسی صرف ایہ نہیں آکھ سکدے کہ $\Nat^2$ صحیح عدداں دا سیٹ اے۔

پچھلی مشکل نوں عام کر کے، ساڈی لوڑ ایہ اے:
	$$a - b = c - d \text{ اُدوں تے صرف اُدوں جدوں }a + d = c + b$$
(ایہ صاف ہونا چاہیدا اے کہ صحیح عدداں نے \emph{ایویں ای ورتاؤ کرنا}
اے: بس دونیں پاسیاں وچ $b$ تے $d$ جمع کرو۔) ایس ورتاؤ دی ضمانت دا
سوکھا طریقہ ایہ اے کہ ترتیب وار جوڑیاں وچکار ہم ارزیت دا تعلق
$\Intequiv$ ایویں تعریف کر دئیے:
$$\tuple{ a, b } \Intequiv \tuple{c, d}\text{ اُدوں تے صرف اُدوں جدوں }a + d = c + b$$
ہُن ساڈے نوں وکھاؤنا اے کہ ایہ ہم ارزیت دا تعلق اے۔
\begin{prop} $\Intequiv$ ہم ارزیت دا تعلق اے۔
\end{prop}
	\begin{proof}
	ساڈے نوں وکھاؤنا اے کہ $\Intequiv$ انعکاسی، متناظر تے متعدی اے۔

	\emph{انعکاسیت:} صاف اے کہ $\tuple{a, b} \Intequiv \tuple{a, b}$،
کیوں جو $a + b = b + a$۔

	\emph{تناظر:} فرض کرو کہ $\tuple{a, b} \Intequiv \tuple{c, d}$، سو
$a + d = c + b$۔ فیر $c + b = a + d$، تے ایس لئی
$\tuple{c, d} \Intequiv \tuple{a, b}$۔

	\emph{تعدیت:} فرض کرو کہ $\tuple{a, b} \Intequiv \tuple{c, d}\Intequiv
\tuple{m, n}$۔ سو $a + d = c + b$ تے $c + n = m + d$۔ ایس لئی
$a + d + c + n = c + b + m + d$، تے فیر $a + n = m + b$۔
لہٰذا $\tuple{a, b } \Intequiv \tuple{m, n}$۔
\end{proof}

ہُن اسی ایس ہم ارزیت دے تعلق نوں ورت کے ہم ارزیت دے طبقے لے سکدے آں:

\begin{defn}
		صحیح عدد، قدرتی عدداں دیاں ترتیب وار جوڑیاں دے $\Intequiv$ تحت
		ہم ارزیت دے طبقے نیں؛ یعنی $\Int = \equivclass{\Nat^2}{\Intequiv}$۔
\end{defn}

ایس مقرر کیتی تعریف بارے کئی وکھ وکھ \emph{فلسفیانہ} ردعمل ہو سکدے نیں۔
پر اوہناں ردعملاں اُتے غور کرن توں پہلاں، کجھ فنی باریکیاں نال اگے
ودھنا فائدے مند اے۔

صحیح عدد کی نیں، ایہ دَس کے، ساڈے نوں اوہناں اُتے بنیادی فنکشن تے تعلق
تعریف کرنے پین گے۔ آؤ $\equivrep{m, n}{\Intequiv}$ اوس ہم ارزیت دے
طبقے لئی لکھئیے جیہڑا $\Intequiv$ تحت اے تے جیہدے وچ
$\tuple{m, n}$، !!a{element} اے۔\footnote{نوٹ: \olref[sfr][rel][eqv]{def:equivalenceclass}
وچ متعارف کرائی علامت ورت کے اسی ایسے شے لئی
$\equivrep{\tuple{m,n}}{\Intequiv}$ لکھدے۔ پر اوہ پڑھنا ذرا اوکھا اے۔}
یعنی:
	$$\equivrep{m, n}{\Intequiv} = \Setabs{\tuple{a, b} \in \Nat^2}{\tuple{a, b}\Intequiv \tuple{m, n}}$$
سو ہُن اسی کجھ تعریفاں پیش کردے آں:
	\begin{align*}
		\equivrep{a, b}{\Intequiv} + \equivrep{c, d}{\Intequiv} &= \equivrep{a + c, b + d}{\Intequiv}\\
		\equivrep{a, b}{\Intequiv} \times \equivrep{c, d}{\Intequiv} &= \equivrep{a c + b  d, a  d + b c}{\Intequiv}\\
		\equivrep{a, b}{\Intequiv} \leq \equivrep{c, d}{\Intequiv} &\text{ اُدوں تے صرف اُدوں جدوں }a + d \leq b + c
	\end{align*}
(عام رواج مطابق، میں ”$ab$“ نوں ”$(a \times b)$“ دا مختصر روپ سمجھ
کے ورت ریا واں، صرف ایس لئی کہ مسلّمے پڑھنے سوکھے رہن۔)
% PNBARI-001 ماخذ دی نحوی درستی: انگریزی عبارت وچ “using ab to stand for” والا to چھڈیا ہویا اے؛ پنجابی صرف فوراً واضح کیتی مراد بیان کردی اے۔
ہُن ساڈے نوں یقین کرنا پینا اے کہ ایہ تعریفاں اوہو جیہا ورتاؤ کردیاں نیں
جیہو جیہا اوہناں نوں \emph{کرنا چاہیدا اے}۔ ایس گل نوں پورا کھولنا تے
جانچنا کافی لمّا کم اے؛ تفصیلاں نوں اسی \olref[check]{sec} ول بھیج دے
آں۔ پر مختصر گل ایہ اے: سب کجھ ٹھیک کم کردا اے!

اک آخری گل رہندی اے۔ اسی قدرتی عدداں نوں ورت کے صحیح عدد بنائے نیں۔
پر ایس دا مطلب ایہ ہووے گا کہ قدرتی عدد \emph{خود صحیح عدد نہیں ہون گے}۔
اسی \olref[ref]{sec} وچ ایس دی فلسفیانہ اہمیت ول واپس آواں گے۔ پر خالص
فنی پاسے توں ساڈے نوں کوئی طریقہ چاہیدا اے جیہدے نال قدرتی عدداں نوں
صحیح عدد \emph{سمجھ کے} ورت سکئیے۔ خیال بڑا سوکھا اے: ہر $n \in \Nat$
لئی اسی بس مقرر کر دیندے آں کہ $n_\Int = \equivrep{n, 0}{\Intequiv}$۔
ساڈے نوں تصدیق کرنی اے کہ ایہ تعریف درست ورتاؤ کردی اے، یعنی ہر
$m, n \in \Nat$ لئی
\begin{align*}
	(m + n)_\Int &= m_\Int + n_\Int\\
	(m \times n)_\Int &= m_\Int \times n_\Int\\
	m \leq n &\liff m_\Int \leq n_\Int
\end{align*}
پر ایہ سب کافی سیدھا اے۔ مثال لئی، ایہ وکھاؤن واسطے کہ ایہناں وچوں
دوجی شرط پوری ہوندی اے، اسی قدرتی عدداں دے ورتاؤ نوں ورت کے ایویں
دلیل دے سکدے آں:
%\begin{proof}
%	قدرتی عدداں دے ورتاؤ نوں ورتدے ہویاں، صاف اے:
	\begin{align*}
%		(m + n)_\Int  = \equivrep{m + n, 0}{\Intequiv} = \equivrep{m + n, 0 + 0}{\Intequiv} = \equivrep{m, 0}{\Intequiv} + \equivrep{n, 0}{\Intequiv} = m_\Int + n_\Int\\
		(m \times n)_\Int  &= \equivrep{m \times n, 0}{\Intequiv} \\
		&= \equivrep{m \times n + 0 \times 0, m \times 0 + 0 \times n}{\Intequiv} \\
		&= \equivrep{m, 0}{\Intequiv} \times \equivrep{n, 0}{\Intequiv} \\
		&= m_\Int \times n_\Int
	\end{align*}
%\end{proof}
باقی دو شرطاں دی تصدیق مشق لئی چھڈدے آں۔
\begin{prob}
وکھاؤ کہ $(m + n)_\Int = m_\Int + n_\Int$ تے
$m \leq n \liff m_\Int \leq n_\Int$، ہر $m, n \in \Nat$ لئی۔
\end{prob}
\end{document}
